\section{让展示适应不同使用条件}
\subsection{布局适配不仅是缩小页面}
在宽屏文章中，图片与文字可以并列；到了手机上，若仍强行保留两列，正文就可能窄得难以阅读。%
因此，本项目通过 Grid、Flexbox 和媒体查询调整布局，文章图文在窄屏转为单列。%
这里的目标是重新安排阅读空间，而不是把整个桌面页面等比例缩小。

进一步调研时，我了解到容器查询：媒体查询常依据视口等条件判断，%
容器查询可以依据组件所在容器的尺寸改变样式\cite{mdn-container}。同一个图文块可能处在宽文章区，%
也可能处在后台狭窄的预览栏，此时容器宽度比“是不是手机”更有意义。其优点是有助于组件复用，%
代价是需要设计查询容器和回退布局。当前项目主要使用媒体查询，容器查询属于我后续研究的方向。

\begin{figure}[H]
\centering
\begin{tikzpicture}[
  font=\fontsize{10.5}{14}\selectfont,
  line width=0.6pt,draw=wmnavy!65,
  labeltext/.style={text=wmnavy,align=left},
  wire/.style={draw=wmnavy!40,line width=0.5pt}
]
\node[anchor=west,font=\fontsize{11}{14}\selectfont\bfseries,text=wmnavy]
  at (0,0.55) {宽容器：图文两列};
\node[anchor=west,font=\fontsize{11}{14}\selectfont\bfseries,text=wmnavy]
  at (10.1,0.55) {窄容器：顺序堆叠};
\draw[rounded corners=3pt] (0,0) rectangle (9.1,-5.0);
\draw[rounded corners=3pt] (10.1,0) rectangle (14.6,-6.25);
\draw[wire] (0,-0.45) -- (9.1,-0.45);
\draw[wire] (10.1,-0.45) -- (14.6,-0.45);
\fill[wmnavy!25] (0.24,-0.23) circle (0.035);
\fill[wmnavy!25] (0.40,-0.23) circle (0.035);
\fill[wmnavy!25] (0.56,-0.23) circle (0.035);
\fill[wmnavy!25] (10.34,-0.23) circle (0.035);
\fill[wmnavy!25] (10.50,-0.23) circle (0.035);
\fill[wmnavy!25] (10.66,-0.23) circle (0.035);
\node[anchor=north west,labeltext,font=\fontsize{11}{14}\selectfont\bfseries]
  at (0.4,-0.85) {01\quad 内容标题};
\node[anchor=north west,labeltext,text width=3.65cm]
  at (0.4,-1.65) {02\quad 正文段落\\以语义顺序组织内容，\\保持段落层级清楚。};
\draw[wire] (0.5,-3.55) -- (3.85,-3.55);
\draw[wire] (0.5,-3.85) -- (3.85,-3.85);
\draw[wire] (0.5,-4.15) -- (3.10,-4.15);
\draw[fill=wmnavy!5,draw=wmnavy!50] (4.75,-1.0) rectangle (8.6,-3.30);
\draw[wire] (4.75,-1.0) -- (8.6,-3.30);
\draw[wire] (8.6,-1.0) -- (4.75,-3.30);
\node[fill=white,text=wmnavy,inner sep=4pt] at (6.675,-2.15) {03\quad 内容图片};
\node[anchor=north west,labeltext,text width=3.7cm] at (4.75,-3.57)
  {04\quad 图注说明};
\node[anchor=north west,labeltext,font=\fontsize{11}{14}\selectfont\bfseries]
  at (10.45,-0.78) {01\quad 内容标题};
\node[anchor=north west,labeltext,text width=3.6cm]
  at (10.45,-1.48) {02\quad 正文段落\\以语义顺序组织内容，\\保持段落层级清楚。};
\draw[fill=wmnavy!5,draw=wmnavy!50] (10.5,-3.42) rectangle (14.2,-5.24);
\draw[wire] (10.5,-3.42) -- (14.2,-5.24);
\draw[wire] (14.2,-3.42) -- (10.5,-5.24);
\node[fill=white,text=wmnavy,inner sep=4pt] at (12.35,-4.33) {03\quad 内容图片};
\node[anchor=north west,labeltext,text width=3.55cm] at (10.45,-5.48)
  {04\quad 图注说明};
\node[anchor=north west,labeltext,text width=8.65cm] at (0,-5.35)
  {同一内容结构：标题 → 正文 → 图片 → 图注。\\视觉列数变化时，阅读与键盘访问顺序保持合理。};
\node[anchor=north west,labeltext,text width=14.4cm] at (0,-6.65)
  {布局响应可用宽度，并尊重“减少动态效果”偏好。图片的 alt 属性提供替代文本，图注补充可见语境；%
    装饰性图片可使用空替代文本。};
\end{tikzpicture}

\caption{图文组件从并列到单列，内容语义保持一致}%
\label{fig:responsive}
\end{figure}

\subsection{动效与可访问性需要一起考虑}
页面入场效果使用浏览器 Web Animations API，并通过 IntersectionObserver
判断元素进入视口的情况\cite{local-code,mdn-waapi}。%
它比为每一个元素单独处理滚动事件更容易组织观察逻辑，但观察回调并不意味着任务没有计算成本，%
动画也不应遮挡交互。项目在减少动画偏好、焦点进入以及能力不足时提供更直接的内容显示方式。

我也逐渐理解，替代文字 alt 与图注并不相同。alt 传达看不到图片时需要知道的信息，%
图注解释画面与正文的联系；一幅装饰图与一张玩法说明图，不应机械使用同样的描述。%
类似地，菜单不能只有鼠标悬停才能使用，状态不能只靠颜色表达，键盘焦点也需要清楚可见。%
WCAG 2.2 提供相应的可访问性要求，但有几个合规属性并不意味着整站已通过审计\cite{wcag22}。

\subsection{用指标衡量体验感}
调研性能时，我接触到 LCP、INP 和 CLS：它们分别关注主要内容显示、交互响应和布局稳定性，%
良好阈值为不超过 2.5 秒、200 毫秒和 0.1，通常按移动端与桌面端分别观察访问的第 75
百分位\cite{web-vitals}。这些指标帮助我把“页面有点慢”拆成更具体的问题，%
但本文没有把行业阈值当作项目实测结果。

例如，图片迟到可能影响主要内容显示；编辑时同步处理很大的正文可能延迟输入反馈；%
没有预留图片空间可能让按钮突然移动。只减小打包文件不能同时解决这些问题，%
单次本地评分也不能代表所有用户。对这个项目，更合理的下一步是固定文章、设备和网络条件，%
分别比较资源请求、关键交互与布局变化，再决定优化重点。
